\documentclass[UTF8,a4paper,12pt]{ctexart}

\usepackage{geometry}
\usepackage{booktabs}
\usepackage{longtable}
\usepackage{array}
\usepackage{hyperref}
\usepackage{xcolor}
\usepackage{listings}
\usepackage{graphicx}
\usepackage{float}

\geometry{left=2.5cm,right=2.5cm,top=2.5cm,bottom=2.5cm}
\hypersetup{colorlinks=true,linkcolor=blue,urlcolor=blue}
\lstset{
  basicstyle=\ttfamily\small,
  breaklines=true,
  frame=single,
  columns=fullflexible,
  backgroundcolor=\color{gray!5}
}

\title{在线书店系统\\大模型驱动的智能化软件测试报告}
\author{第二部分负责人：使用 Codex 协作完成}
\date{2026年4月24日}

\begin{document}
\maketitle
\tableofcontents
\newpage

\section{项目与作业目标}

本报告面向软件测试课程作业中“\textbf{大模型驱动的智能化软件测试}”部分，选取课程项目“在线书店系统”作为被测对象，使用 Codex 作为软件开发 Agent，完成需求分析、测试计划、黑盒测试设计、黑盒自动化执行、白盒测试结果整理和测试报告撰写。

本报告遵循以下原则：
\begin{itemize}
  \item 所有黑盒测试通过数、失败数、执行命令、问题修复结论均来自真实执行。
  \item 黑盒自动化执行口径统一为 Playwright 浏览器端到端测试。
  \item 白盒部分保留现有实现与 JaCoCo 统计结果，不在本轮黑盒调整中变更。
\end{itemize}

\section{被测项目说明}

\subsection{项目基本信息}

\begin{longtable}{p{4cm}p{10cm}}
\toprule
项目名称 & 在线书店系统 \\
\midrule
前端技术栈 & React、React Router、Ant Design、fetch API \\
后端技术栈 & Spring Boot、Spring Data JPA、MySQL、Session 认证 \\
需求说明文档 & \texttt{README.md} \\
前端源码目录 & \texttt{src/} \\
后端源码目录 & \texttt{backend/src/main/java/} \\
测试文档目录 & \texttt{docs/testing/} \\
\bottomrule
\end{longtable}

\subsection{源码规模}

本次统计只包含真正的前后端业务代码，不包含 \texttt{public}、Markdown 文档、测试代码和生成目录。

\begin{longtable}{p{6cm}p{6cm}}
\toprule
统计项 & 结果 \\
\midrule
后端代码文件数 & 52 \\
后端代码总行数 & 3694 行 \\
前端代码文件数 & 26 \\
前端代码总行数 & 4008 行 \\
前后端总代码文件数 & 78 \\
前后端总代码总行数 & 7702 行 \\
\bottomrule
\end{longtable}

因此，本项目满足课程要求中“源码代码行数不低于 1000 行”的约束。

\section{需求功能点提取}

根据 \texttt{README.md}，本项目至少包含以下 16 个核心可测试功能点：

\begin{longtable}{p{1.5cm}p{8cm}p{3cm}}
\toprule
编号 & 功能点 & 角色 \\
\midrule
F01 & 用户注册 & 普通用户 \\
F02 & 用户登录 & 普通用户、管理员 \\
F03 & 登录状态校验 & 全部角色 \\
F04 & 未登录访问受保护页面拦截 & 全部角色 \\
F05 & 图书列表浏览 & 全部角色 \\
F06 & 图书搜索 & 全部角色 \\
F07 & 图书详情查看 & 全部角色 \\
F08 & 加入购物车 & 普通用户 \\
F09 & 购物车查看 & 普通用户 \\
F10 & 购物车结算 & 普通用户 \\
F11 & 订单查询 & 普通用户 \\
F12 & 管理员订单管理 & 管理员 \\
F13 & 管理员图书管理 & 管理员 \\
F14 & 图书销量统计 & 管理员 \\
F15 & 用户消费统计 & 管理员 \\
F16 & 个人购书统计 & 普通用户 \\
\bottomrule
\end{longtable}

普通用户核心流程定义为：

\begin{center}
\texttt{登录 -> 浏览图书 -> 进入详情 -> 加入购物车 -> 勾选商品结算 -> 查询订单}
\end{center}

管理员核心流程定义为：

\begin{center}
\texttt{管理员登录 -> 订单管理 -> 图书管理 -> 图书销量统计 -> 用户消费统计}
\end{center}

\section{Codex 驱动的测试工作流}

\subsection{工作流阶段}

\begin{longtable}{p{2cm}p{5cm}p{7cm}}
\toprule
阶段 & 输入 & 输出 \\
\midrule
需求分析 & \texttt{README.md} & 功能点清单、核心流程、测试范围 \\
黑盒设计 & 需求与页面流程 & 等价类、边界值、决策表用例 \\
黑盒实现 & 黑盒设计文档 & Playwright 浏览器端到端黑盒代码 \\
黑盒执行 & 前后端运行环境 & 真实通过数、失败数、问题修复记录 \\
白盒整理 & 既有测试与 JaCoCo 产物 & 覆盖率和方法说明 \\
结果沉淀 & 执行产物 & 测试报告、映射表、LaTeX 文档 \\
\bottomrule
\end{longtable}

\subsection{当前黑盒自动化代码}

\begin{itemize}
  \item \texttt{playwright.config.js}
  \item \texttt{e2e/decision-boundary.spec.js}
  \item \texttt{e2e/user-purchase-flow.spec.js}
  \item \texttt{e2e/admin-statistics.spec.js}
\end{itemize}

说明：当前仓库中的黑盒自动化已经完全统一为 Playwright，不再保留后端 API 级黑盒或前端页面级黑盒执行层。

% 图1：建议插入 Codex 阅读 README 并生成黑盒方案的对话截图。
\begin{figure}[H]
  \centering
  % \includegraphics[width=0.9\textwidth]{figures/codex-blackbox-plan.png}
  \caption{Codex 生成黑盒测试方案的交互截图（示意位置）}
\end{figure}

\section{测试计划}

\subsection{测试目标}

\begin{itemize}
  \item 验证在线书店核心业务链路是否正确。
  \item 验证普通用户和管理员的权限隔离是否正确。
  \item 通过 Playwright 将黑盒测试落到真实浏览器环境中执行。
  \item 保留现有白盒测试结果，满足课程对白盒覆盖率的要求。
\end{itemize}

\subsection{测试环境}

\begin{longtable}{p{4cm}p{10cm}}
\toprule
类别 & 内容 \\
\midrule
操作系统 & Windows，项目路径为 \texttt{D:\textbackslash course\textbackslash softwaretest\textbackslash hw1\textbackslash web-ebookstore} \\
黑盒执行框架 & Playwright \\
黑盒前端地址 & \texttt{http://localhost:3000} \\
黑盒后端地址 & \texttt{http://localhost:8081} \\
黑盒测试账号 & 管理员：\texttt{admin / 123456}；普通用户：\texttt{coco / 123456} \\
白盒覆盖率工具 & JaCoCo 0.8.12 \\
\bottomrule
\end{longtable}

\section{黑盒测试设计}

\subsection{方法说明}

黑盒测试基于需求说明文档设计，不依赖内部实现，采用如下三类方法：

\begin{itemize}
  \item 等价类：验证登录成功、登录失败、搜索空结果、订单查询等典型输入类别。
  \item 边界值：验证注册用户名最小长度、确认密码一致性、购物车未勾选商品等临界场景。
  \item 决策表：验证“是否登录”“是否管理员”“是否勾选商品”等条件组合对系统行为的影响。
\end{itemize}

\subsection{代表性黑盒用例}

\begin{longtable}{p{1.8cm}p{3.5cm}p{2.3cm}p{4cm}p{4cm}}
\toprule
编号 & 功能点 & 方法 & 输入/操作 & 预期结果 \\
\midrule
BB-02 & 注册用户名过短 & 边界值 & 用户名长度小于 3 & 前端校验拦截 \\
BB-03 & 注册确认密码不一致 & 边界值 & 两次输入密码不同 & 前端校验拦截 \\
BB-04 & 普通用户登录成功 & 等价类 & 输入 \texttt{coco / 123456} & 登录成功 \\
BB-05 & 密码错误登录失败 & 等价类 & 正确用户名 + 错误密码 & 登录失败 \\
BB-06 & 未登录访问购物车 & 决策表 & 直接访问 \texttt{/cart} & 重定向到登录页 \\
BB-07 & 普通用户访问管理员页面 & 决策表 & 普通用户访问 \texttt{/admin/orders} & 显示权限不足 \\
BB-11 & 搜索不存在的图书关键字 & 等价类 & 输入不存在关键字 & 显示空结果提示 \\
BB-13 & 未勾选商品不可结算 & 边界值 & 购物车有商品但不勾选 & 结算按钮禁用 \\
BB-14 & 勾选商品后结算成功 & 决策表 & 已登录 + 已勾选商品 & 成功生成订单 \\
BB-15 & 按书名查询订单 & 等价类 & 在订单页输入书名搜索 & 返回匹配订单 \\
BB-16 & 管理员订单管理 & 决策表 & 管理员进入订单管理页 & 加载全部订单 \\
BB-17 & 管理员图书管理 & 决策表 & 管理员进入图书管理页 & 加载图书表格 \\
BB-18 & 图书销量统计 & 决策表 & 管理员发起销量查询 & 返回销量统计表格 \\
BB-19 & 用户消费统计 & 决策表 & 管理员切换页签并查询 & 返回用户消费统计 \\
BB-20 & 个人购书统计 & 等价类 & 普通用户进入个人统计页 & 返回个人购书统计结果 \\
\bottomrule
\end{longtable}

\subsection{典型决策表}

\subsubsection{登录决策表}

\begin{longtable}{p{2cm}p{2.5cm}p{2.5cm}p{3cm}p{4cm}}
\toprule
规则 & 用户存在 & 密码正确 & 权限/状态满足 & 预期结果 \\
\midrule
R1 & 否 & 任意 & 任意 & 登录失败 \\
R2 & 是 & 否 & 任意 & 登录失败 \\
R3 & 是 & 是 & 否 & 登录失败或拒绝访问 \\
R4 & 是 & 是 & 是 & 登录成功 \\
\bottomrule
\end{longtable}

\subsubsection{购物车结算决策表}

\begin{longtable}{p{2cm}p{2.5cm}p{2.5cm}p{2.5cm}p{4cm}}
\toprule
规则 & 已登录 & 已选中商品 & 库存充足 & 预期结果 \\
\midrule
R1 & 否 & 任意 & 任意 & 跳转登录页 \\
R2 & 是 & 否 & 任意 & 按钮禁用，不允许结算 \\
R3 & 是 & 是 & 否 & 结算失败 \\
R4 & 是 & 是 & 是 & 成功生成订单 \\
\bottomrule
\end{longtable}

\section{Playwright 黑盒自动化真实执行结果}

\subsection{真实执行方式}

先手动启动前端：

\begin{lstlisting}[language=bash]
cd D:\course\softwaretest\hw1\web-ebookstore
npm start
\end{lstlisting}

再手动启动后端：

\begin{lstlisting}[language=bash]
cd D:\course\softwaretest\hw1\web-ebookstore\backend
mvn spring-boot:run
\end{lstlisting}

真实执行黑盒命令：

\begin{lstlisting}[language=bash]
cmd /c "set PLAYWRIGHT_MANUAL_SERVER=1 && npx playwright test --workers=1"
\end{lstlisting}

\subsection{真实执行结果}

2026年4月24日真实执行结果如下：

\begin{longtable}{p{6cm}p{6cm}}
\toprule
指标 & 结果 \\
\midrule
测试文件数量 & 3 \\
浏览器端到端测试用例总数 & 13 \\
失败数 & 0 \\
错误数 & 0 \\
执行结果 & PASS \\
总耗时 & 38.0s \\
\bottomrule
\end{longtable}

按文件划分的真实结果为：

\begin{itemize}
  \item \texttt{e2e/decision-boundary.spec.js}：6 passed
  \item \texttt{e2e/user-purchase-flow.spec.js}：2 passed
  \item \texttt{e2e/admin-statistics.spec.js}：5 passed
\end{itemize}

% 图2：建议插入 Playwright 最终 13 passed 的执行截图。
\begin{figure}[H]
  \centering
  % \includegraphics[width=0.9\textwidth]{figures/playwright-13-passed.png}
  \caption{Playwright 浏览器端到端黑盒自动化执行结果截图（示意位置）}
\end{figure}

\subsection{当前已覆盖的黑盒内容}

\begin{itemize}
  \item 未登录访问购物车时重定向到登录页
  \item 用户名正确但密码错误时登录失败
  \item 普通用户访问管理员订单页时显示权限不足
  \item 注册用户名长度小于 3 时被前端校验拦截
  \item 注册确认密码不一致时被前端校验拦截
  \item 购物车未勾选任何商品时结算按钮禁用
  \item 搜索不存在的图书关键字时显示空结果提示
  \item 普通用户登录成功
  \item 浏览图书列表
  \item 进入图书详情页
  \item 将图书加入购物车
  \item 勾选商品后成功结算
  \item 订单页按书名查询订单
  \item 管理员订单管理
  \item 管理员图书管理
  \item 图书销量统计
  \item 用户消费统计
  \item 个人购书统计
  \item 跑通“登录 -> 浏览 -> 加购 -> 下单 -> 查单”核心业务链路
\end{itemize}

\subsection{真实发现并修复的问题}

Playwright 第一轮全量真实执行时，13 个用例全部失败。失败原因不是业务逻辑错误，而是黑盒脚本对登录页和注册页的定位器假设错误：

\begin{itemize}
  \item 错误假设页面存在 \texttt{form[name="login"]} 结构。
  \item 登录与注册页签中存在同名输入框，触发了 Playwright 严格模式冲突。
\end{itemize}

修复方式：

\begin{itemize}
  \item 改为基于真实可见控件定位用户名和密码输入框。
  \item 登录按钮改为定位激活页签中的主按钮。
  \item 注册页输入框改为在激活页签中按输入顺序定位。
\end{itemize}

修复后重新真实回归，最终全量结果为 13 passed，说明当前 Playwright 黑盒方案已经稳定可执行。

\section{白盒测试结果说明}

本轮调整只涉及黑盒自动化，既有白盒测试实现和 JaCoCo 统计结果保持不变。白盒覆盖率数据沿用既有真实执行产物：

\begin{itemize}
  \item \texttt{backend/target/site/jacoco/jacoco.csv}
\end{itemize}

当前覆盖率结果如下：

\begin{longtable}{p{6cm}p{6cm}}
\toprule
覆盖指标 & 结果 \\
\midrule
总行覆盖率 & 97.97\% \\
总指令覆盖率 & 96.46\% \\
总分支覆盖率 & 82.54\% \\
\bottomrule
\end{longtable}

这部分满足课程对白盒测试覆盖率不低于 95\% 的要求。

\section{未来功能变更适配方案}

\subsection{新增用户角色}

可能变更：系统新增“运营人员”或“商家”角色。

\begin{itemize}
  \item 黑盒调整：扩展权限决策表，补充新角色访问订单、图书和统计页面的 Playwright 场景。
  \item 白盒调整：补充拦截器和控制器角色判断分支覆盖。
\end{itemize}

\subsection{支付流程细化}

可能变更：订单状态扩展为“待支付、已支付、已取消”。

\begin{itemize}
  \item 黑盒调整：增加支付成功、支付失败、超时取消等决策表场景。
  \item 白盒调整：补充订单状态机的基本路径覆盖和数据流分析。
\end{itemize}

\subsection{新增优惠券功能}

可能变更：下单时支持优惠券。

\begin{itemize}
  \item 黑盒调整：增加有效券、过期券、已使用券、金额边界等 Playwright 场景。
  \item 白盒调整：补充订单金额计算与优惠券状态更新的数据流测试。
\end{itemize}

\section{课程要求逐条对应}

\begin{longtable}{p{6cm}p{3cm}p{6cm}}
\toprule
课程要求 & 是否满足 & 说明 \\
\midrule
源码代码行数不少于 1000 行 & 满足 & 前后端生产代码共 7702 行 \\
功能点不少于 10 个 & 满足 & 本文列出 16 个核心功能点 \\
至少包含一个核心流程 & 满足 & 覆盖普通用户购买链路和管理员管理链路 \\
包含需求说明文档 & 满足 & 使用 \texttt{README.md} 作为需求依据 \\
使用至少一种软件开发 Agent & 满足 & 使用 Codex \\
完成测试计划 & 满足 & 本文第 5 节及相关文档已覆盖 \\
黑盒包含边界值、等价类、决策表 & 满足 & 本文第 6 节和第 7 节已覆盖 \\
白盒包含代码覆盖、基本路径覆盖、数据流分析 & 满足 & 既有白盒文档和 JaCoCo 结果已覆盖 \\
白盒覆盖率不低于 95\% & 满足 & 总行覆盖率 97.97\% \\
完成测试报告 & 满足 & 本文和 \texttt{docs/testing/test-report.md} 已形成可提交成果 \\
\bottomrule
\end{longtable}

\section{小组成员贡献}

\begin{longtable}{p{4cm}p{8cm}p{3cm}}
\toprule
成员 & 贡献内容 & 比例 \\
\midrule
第二部分负责人 & 使用 Codex 完成需求分析、黑盒设计、Playwright 黑盒自动化、真实执行、问题修复、白盒结果整理与报告撰写 & 100\% \\
\bottomrule
\end{longtable}

\section{提交材料清单}

\begin{itemize}
  \item 被测程序源码：整个 \texttt{web-ebookstore} 项目
  \item 需求说明文档：\texttt{README.md}
  \item 测试计划：\texttt{docs/testing/test-plan.md}
  \item 黑盒测试文档：\texttt{docs/testing/test-cases-blackbox.md}
  \item 白盒测试文档：\texttt{docs/testing/test-cases-whitebox.md}
  \item 黑盒映射表：\texttt{docs/testing/blackbox-automation-mapping.md}
  \item 测试报告：\texttt{docs/testing/test-report.md}
  \item 大模型使用记录：\texttt{docs/testing/codex-workflow.md}
  \item 覆盖率报告：\texttt{backend/target/site/jacoco/index.html}
\end{itemize}

% 图3：建议插入“基于大模型的软件测试使用记录”截图。
\begin{figure}[H]
  \centering
  % \includegraphics[width=0.9\textwidth]{figures/llm-usage-record.png}
  \caption{大模型驱动测试工作流使用记录截图（示意位置）}
\end{figure}

\section{结论}

当前仓库中的黑盒自动化已经完全统一为 Playwright 浏览器端到端方案。2026 年 4 月 24 日，在手动启动前端和后端后，真实执行
\texttt{cmd /c "set PLAYWRIGHT\_MANUAL\_SERVER=1 \&\& npx playwright test --workers=1"}，
共得到 13 个测试全部通过。

这些场景覆盖了等价类、边界值和决策表三类黑盒设计方法，且真实跑通了在线书店的用户核心购买链路和管理员核心管理链路。白盒部分继续保持既有实现与 JaCoCo 覆盖率结果，总行覆盖率为 97.97\%，满足课程对白盒覆盖率不低于 95\% 的要求。

因此，当前成果已经形成一套可提交课程作业的“大模型驱动智能化软件测试”方案，黑盒部分具有统一执行口径、真实执行结果和可复现能力。

\end{document}
